\documentclass[10pt,a4paper]{amsart}
\usepackage[margin=19mm]{geometry}
\usepackage{amsmath,amssymb,mathtools}
\usepackage[scr=rsfs,cal=euler]{mathalfa}
\usepackage{xfrac}
\usepackage[unicode,hidelinks,bookmarks=false]{hyperref}
\newcommand{\ie}{i.e.\ }
\newcommand{\cf}{cf.\ }
\newcommand{\resp}{respectively\ }
\newcommand{\Subsection}{\subsection}
\providecommand{\nobreakdash}{\nobreak\hbox{-}}
\setlength{\emergencystretch}{2em}
\multlinegap0pt
\usepackage{bm}
\usepackage{bbm}
\usepackage{dsfont}
\usepackage{xspace}
\usepackage{cancel}
\usepackage{mathtools}
\usepackage{amsthm}
\makeatletter
\@ifundefined{thm}{\newtheorem{thm}{Theorem}}{}
\makeatother
\newcommand{\IR}{\text{${\mathbb{R}}$}}
\newcommand{\IZ}{\text{${\mathbb{Z}}$}}
\newcommand{\del}{\mbox{$\partial$}}
\newcommand{\half}{\mbox{$\frac{1}{2}$}}
\newcommand{\p}{\mathfrak{p}}
\newcommand{\q}{\mathfrak{q}}
\newcommand{\s}{\mathfrak{s}}
\title[On the Hartree-Fock Eigenvalue Problem]{On the Hartree-Fock Eigenvalue Problem}
\author{Richard A Zalik}
\date{Source: arXiv:2408.07708v3 (CC BY 4.0)}
\begin{document}
\maketitle

\noindent\emph{Source and licence.} On the Hartree-Fock Eigenvalue Problem. Version arXiv:2408.07708v3 (CC BY 4.0), distributed under the Creative Commons Attribution 4.0 International licence, as recorded in the arXiv licence metadata for this exact version and shown on the abstract page https://arxiv.org/abs/2408.07708v3. Full rights notice: licenses/arxiv-2408.07708-CC-BY-4.0.txt.\par\smallskip

\noindent\textit{Editorial scope.} Counted unit: the Thm whose source label is \texttt{thm4} in the article (source0.tex lines 441--458, source label \texttt{thm4}), delivered together with the article's own complete original proof of it (source0.tex lines 459--490, one \texttt{proof} environment of 32 lines). No mathematics is written, repaired or re-derived: statement and proof are the article's own text line for line. Required context retained uncounted: thm1 (lines 197--199); eq4 (lines 227--233); eq5 (lines 256--258); eq6 (lines 303--305); thm3 (lines 411--422). Every definition, display and label of those lines is retained unchanged. Omitted from this package and not redistributed: 1--196, 200--226, 234--255, 259--302, 306--410, 423--440, 491--585. Nothing inside a retained excerpt is omitted or altered: every retained body is delivered exactly as the article's lines are written, so no omission receipt of any kind accompanies this package. Citations: the retained excerpts carry no citation command at all, so no bibliography entry of the article is transcribed and no reference is redistributed. Reference adaptation: none was needed and none was made; every reference inside the delivered unit resolves inside it, so no unresolved reference or citation is produced. Printed slips kept literal: no printed word, symbol, hypothesis or number of the counted statement or of its proof was altered. Layout and class: the article's own class is \texttt{amsart} with packages including babel, setspace, mathtools, amsmath, amsthm, amssymb, amsfonts, graphics; the wrapper is the lane's pinned \texttt{amsart} editorial wrapper at 10pt on A4, which restates the theorem-like environments the retained text needs and adds the article's own macro definitions that the retained text needs (\texttt{\textbackslash IR}, \texttt{\textbackslash IZ}, \texttt{\textbackslash del}, \texttt{\textbackslash half}, \texttt{\textbackslash p}, \texttt{\textbackslash q}, \texttt{\textbackslash s}), with extra wrapper packages (bm, bbm, dsfont, xspace, cancel, mathtools). The delivered numbering is the unit's own. Assets: no retained span contains a figure environment, a graphics-inclusion command, a PDF-page inclusion, a table or a TikZ picture; no image file is copied into this package, the package declares no asset dependency and no image file is redistributed with it. Licence: version arXiv:2408.07708v3 of this article is distributed under the Creative Commons Attribution 4.0 International licence, as recorded in the arXiv licence metadata for this exact version and shown on the abstract page https://arxiv.org/abs/2408.07708v3. A full rights notice is delivered at licenses/arxiv-2408.07708-CC-BY-4.0.txt. Characters: every retained excerpt is pure ASCII, so the delivered engine, XeLaTeX with its default Latin Modern text font, reports no missing character.
\begin{thm} \label{thm1} Let $V$ denote the set of points in $\IR^3$ other than the singular points $\xi_c$; then 
$\nabla^2\psi_a \in C(V) \cap L_2$. If $\psi_a$ is also in $L_1$ (or if $\lim_{x\rightarrow \infty}\psi_a(x)=0$), then $\nabla^2\psi_a $ is in $L_1$ as well.
\end{thm}

\begin{multline} \label{eq4}
||h_{\tau}\psi_a||^2_2=\int |h_{\tau}(s)\psi_a(s)|^2 \,ds= \\
\int |s-\tau|^{-2}\,|\psi_a(s)|^2 \,ds=\int |s|^{-2}\, |\psi_a(s+\tau)|^2 \,ds =\\
\int _{|s|\le1}|s|^{-2}|\psi_a(s+\tau)|^2 \,ds + \int_{|s| >1} |s|^{-2}|\psi_a(s+\tau)|^2 \,ds\le\\
\left (\int _{|s|\le1}|s|^{-2}\,ds \right)^{\half} \left (\int _{|s|\le1}|\psi_a(s+\tau)|^2 \,ds \right )^{\half} +\int_{|s| >1}|\psi_a(s+\tau)|^2 \,ds \le \\
4\pi+1.
\end{multline}

\begin{equation} \label{eq5}
|\s_{a,c}| \le 2\sqrt{\pi}+1.
\end{equation}

\begin{equation} \label{eq6}
\del_j\psi(x), \quad \del_jw(x), \quad \del^2_j\psi(x),\; \text{and}\quad\del^2_j w(x), \quad j=1, 2, 3,
\end{equation}

\begin{thm} \label{thm3}
Let $w$ be such that its partial derivatives of order less or equal to 2 displayed on \eqref{eq6} exist almost everywhere, that $w$ and its partial derivatives of order 2  displayed on \eqref{eq6} are in $ L_1\cap L_2 \cap L_{\infty}$,  and that
\[
\nabla^2[\psi_a \ast w](x)= [\psi_a\ast \nabla^2w](x).
\]
Then for $x \in \IR^3$ and $a=1, \dots n$,
\begin{equation} \label{eq10}
 \varepsilon_a [\psi_a \ast w](x)=
\left [\left ( (\q - \p) \psi _a-2\sum_{c=1}^n \s_{a,c}\psi_{c}\right )\ast w\right ](x) -[\psi_{a}\ast \nabla^2w](x)
\end{equation}
and $\psi _a\ast w$ is in $C_0$  and is uniformly continuous on $\IR^3$.
\end{thm}

\begin{thm}\label{thm4}\hspace{1in}\\
(a)\ Let
\begin{equation} \label{eq11}
 \varepsilon_{a,\nu}= 
 \frac{\big{\Vert} [(\q -\p )T_{a,\nu } -2\sum_{c=1}^n(\s_{a,c,\nu} T_{c,\nu}) ] \ast w -  [ T_{a, \nu}\ast \nabla^2w ] \big{\Vert}_2}
 {||T_{\nu,a} \ast w ||_2};
 \end{equation}
then $\lim_{\nu\rightarrow \infty}\varepsilon_{a,\nu}=\varepsilon_a$.\\
(b)\ $\{[T_{\nu,a}\ast w](x); \nu \in \IR^3\}$ converges to $[\psi_a\ast w](x)$ uniformly on $\IR^3$, and
\begin{equation} \label{eq12}
|| \psi_a\ast w-T_{\nu,a}\ast w||_u \le |||\psi_a-T_{\nu,a}||_2\,||w||_2.
\end{equation}
(c)\ Assume that $[\psi_a\ast w](x_0)\ne 0$; then for $\nu$ sufficiently large $[T_{a,\nu}\ast w](x_0) \ne 0$ and, if
\begin{equation} \label{eq13}
\varepsilon_{\nu,a}(x_0)= \frac{\{[(\q _{\nu}-\p )T_{a,\nu } -2\sum_{c=1}^n(\s_{a,c,\nu} T_{c,\nu})\ast] w \}(x_0)-  [ T_{a, \nu}\ast \nabla^2w ] (x_0)}{[T_{a,\nu}\ast w](x_0)},
\end{equation}
then $\lim_{\nu \rightarrow \infty}\varepsilon_{a,\nu}(x_0)= \varepsilon_a$.
\end{thm}

\begin{proof}
(a)\ Since $\psi -  T_{\nu,a} \in L_2$, Young's inequality implies that $||T_{\nu,a} \ast w ||_2$ converges to $||\psi \ast w ||_2$ and that
the  numerator of \eqref{eq11} converges to 
\[
||[(\q\psi_{a} - \p \psi_{a} -2\sum_{c=1}^n (\s_{a,c}\psi_{c})\ast w ] -[\psi_{a}\ast \nabla^2w]||_2,
\]
and the assertion follows from Theorem \ref{thm3}.\\
(b)\ Inequality \eqref{eq12} follows from Young's inequality.\\
(c)\ That the denominator of \eqref{eq13} does not vanish for sufficiently large $\nu$, as well as the convergence of $(T_{a, \nu} \ast w)(x_0)$ to
$(\psi_a \ast w)(x_0)$, follow from \eqref{eq12}. 
Proceeding as in the proof of Theorem \ref{thm1} we see that
\begin{multline*}
||h_{\tau}T_{a,\nu} ||^2_2=\int |h_{\tau}(s)T_{a,\nu}(s)|^2 \,ds \le\\
\left (\int _{|s|\le1}|s|^{-2}\,ds \right)^{\half} \left (\int _{|s|\le1}|T_{a,\nu}(s+\tau)|^2 \,ds \right )^{\half} +\int_{|s| >1}|T_{a,\nu}(s+\tau)|^2 \,ds \le\\
4\pi K_a+K^2_a,
\end{multline*}
which implies that $||\p T_{a,\nu}||_2$ is bounded, uniformly on $\nu$, and by another application of Young's inequality
as in \eqref{eq12} above we conclude that $[(\p T_{a,\nu})\ast w](x_0)$ converges to $[(\p\psi_a)\ast w](x_0)$.
From \eqref{eq12} we know there is a constant $M_a$ such that $||T_{a,\nu}||_u \le M_a$. Repeating the arguments used in the proof of Theorem 1 we readily obtain
\begin{multline*}
|\s_{a,c,\nu}(x)| \le\\
 \left ( \int_{|s|\le 1}|T^2_{a,\nu}(x-s) T^2_{c,\nu}(x-s)|\,ds \right )^{\half} \left (\int_{|s|\le1}|s|^{-2}\, ds \right)^{\half} +\\
\int_{|s| > 1}|T_{a,\nu}(x-s) T_{c,\nu}(x-s)| \,|s|^{-1}\,ds\le\\
2 \sqrt{\pi} M_a^2+K^2_a,
\end{multline*}
which also implies the boundedness of $\{\q_{\nu}; \nu \in \IZ^3\}$ in $L_{\infty}$. The uniform bounds for 
$||h_{\tau}T_{a,\nu }||_2$ and $\s_{a,c,\nu}$ that we have just obtained are the counterparts of \eqref{eq4} and \eqref{eq5}, and using them we can readily show that the sequence 
\[
\{[( \q_{\nu}-\p )T_{a,\nu} -2\sum_{c=1}^n(\s_{a,c,\nu} T_{c,\nu}) ); \nu \in \IZ^3\}
\]
is bounded in $L_2$, and the assertion follows by convolving with $w$ and applying Young's inequality as in \eqref{eq12}.
\end{proof}

\end{document}
